\section{Experimental design}
\label{sec:design}

% STAGE 3: rewritten 8 September 2026.  Table~\ref{tab:arms} is new: one row
% per experimental arm, so the reader has a map of every run in the paper.

\subsection{Factors and arms}

The design varies five factors. \textbf{Architecture}: FEDOT.LLM, AutoDS-Tools,
Terminus-2. \Ktool{} and \Kdisc{}: each on or off, crossed on AutoDS-Tools and Terminus-2.
Terminus-2 ships no prescription of its own, so the halves of the AutoDS-Tools
layer are appended to its task file. \textbf{Backbone}:
\texttt{gemma-4-31b-it} throughout. AutoDS-Tools and Terminus-2 are also run
on a smaller model of the same family, \texttt{gemma-4-26b-a4b-it}, and on a
larger one from another family, \texttt{glm-4.7}, because the Gemma family has
nothing larger than 31B, so the upper point differs in family as well as in
size. Starting point: greenfield, where the agent receives the task statement alone,
and improve-a-baseline, where it also receives a working training script. Provisioning, the container image, is varied once, on Terminus-2. Swapping the
stock image for the enriched one changes the installed libraries and leaves the
instruction unchanged. Table~\ref{tab:arms} lists every arm these factors produce, with
the number of trials and the section that reports it.

\begin{table*}[t]
\centering
\caption{Every experimental arm in the paper. A trial is one attempt at one task; TabReD arms cover eight tasks, MLAgentBench arms the ten tasks of the benchmark, of which eight could be run three times on our hardware (Section~\ref{sec:coverage}). Stock: the benchmark's own image; enriched: the image built for AutoDS-Tools, with the libraries its layer names. Welded in: the library choice is fixed by the architecture (Section~\ref{sec:axis}); trimmed: the discipline half without its two clauses about training time (Section~\ref{sec:kdiscterm}). FEDOT image: the image built for FEDOT.LLM (Section~\ref{sec:deviations}).}
\label{tab:arms}
\footnotesize
\setlength{\tabcolsep}{4pt}
\begin{tabular}{lllllrl}
\toprule
System & Benchmark & Layer & Image & Backbone & Trials & Reported in \\
\midrule
FEDOT.LLM    & TabReD & welded in & FEDOT image & gemma-4-31b & 24 & \S\ref{sec:halfrange} \\
Terminus-2   & TabReD & none & stock & gemma-4-31b & 24 & \S\ref{sec:halfrange} \\
Terminus-2   & TabReD & none & enriched & gemma-4-31b & 24 & \S\ref{sec:matched} \\
Terminus-2   & TabReD & none, \Kdisc{}, \Ktool{}, both, \Kdisc{} trimmed & enriched & gemma-4-31b & $5\times24$ & \S\ref{sec:kdiscterm} \\
Terminus-2   & TabReD & none & enriched & gemma-4-26b-a4b, glm-4.7 & $2\times24$ & \S\ref{sec:modelaxis} \\
AutoDS-Tools & TabReD & none, \Kdisc{}, \Ktool{}, both & enriched & gemma-4-31b & $4\times24$ & \S\ref{sec:halfrange}, \S\ref{sec:grid} \\
AutoDS-Tools & TabReD & both & enriched & gemma-4-26b-a4b, glm-4.7 & $24+8$ & \S\ref{sec:modelaxis} \\
\midrule
AutoDS-Tools & MLAgentBench & none & enriched & gemma-4-31b & 30 & \S\ref{sec:headtohead}, \S\ref{sec:signflip} \\
AutoDS-Tools & MLAgentBench & both & enriched & gemma-4-31b & 30 & \S\ref{sec:signflip} \\
Terminus-2   & MLAgentBench & none & stock & gemma-4-31b & 24 & \S\ref{sec:headtohead} \\
FEDOT.LLM    & MLAgentBench & welded in & stock & gemma-4-31b & 2 & \S\ref{sec:coverage} \\
\midrule
AutoDS-Tools & three case studies & maximal, hand-written & enriched & gemma-4-31b & 3 & \S\ref{sec:cases} \\
Terminus-2   & three case studies & none, maximal & enriched & gemma-4-31b & $2\times9$ & \S\ref{sec:cases} \\
FEDOT.LLM    & two tabular cases & welded in & FEDOT image & gemma-4-31b & $2\times3$ & \S\ref{sec:cases} \\
\bottomrule
\end{tabular}
\end{table*}

\subsection{Benchmarks and reference points}
\label{sec:benchmarks}

\textbf{TabReD} contributes eight real-world tabular tasks with temporal
splits~\citep{rubachev2024tabred} and a published reference table of eighteen
methods: gradient boosting implementations~\citep{chen2016xgboost,ke2017lightgbm,prokhorenkova2018catboost},
tabular deep learning architectures, ensembles and retrieval-augmented models,
each tuned with Optuna~\citep{akiba2019optuna} on the official validation split
and averaged over fifteen seeds. Metrics are ROC-AUC for classification and
RMSE in raw target units for regression. The benchmark ships no agentic
protocol, so tasks are solved from scratch: this is the greenfield condition.
Every TabReD result in the paper is read against the reference table in two
ways: as an average rank among the published methods and the agents, and on a
normalised scale. For a task $j$ with score $s$ the normalised score is
$n_j(s) = (s - w_j)/(b_j - w_j)$, where $w_j$ and $b_j$ are the weakest and the
strongest published values on that task, so that $0$ is the weakest and $1$ the
strongest published method; for RMSE the weakest value is the largest. The mean
over the eight tasks has no floor, so a system below the weakest published
method on one task can score below $0$ there.

\textbf{MLAgentBench} contributes ten tasks spanning tabular, text, vision and
graph data~\citep{huang2024mlagentbench}. Each supplies a working but weak
training script and asks the agent to improve it: this is the
improve-a-baseline condition. MLAgentBench publishes no reference table, and the literature offers none at
the granularity we would need. MLAgentBench scores therefore stay in each
task's own metric, with no external ceiling to normalise against. Every MLAgentBench number in the paper compares two of our own runs under
identical conditions, for example AutoDS-Tools with the layer on against the
same system with it off, or AutoDS-Tools against Terminus-2 on the same task.
The normalised scale of Section~\ref{sec:ceiling} is used for TabReD only.

\textbf{Three scientific case studies}, published studies with author-reported
baselines, test external validity outside benchmark conditions
(Section~\ref{sec:cases}). Their instructions prescribe more than any other condition in the paper: the
library, the featurisation, the leakage handling and the baseline to beat. They
are therefore the most heavily prescribed point on the axis of
Section~\ref{sec:axis}. Terminus-2 runs on the full instruction and on the same instruction with the
library section removed. The case studies thus cross tool prescription with
architecture a second time, after TabReD (Sections~\ref{sec:grid}
and~\ref{sec:kdiscterm}).

\subsection{Execution environment and provisioning}

All runs execute inside the Harbor framework~\citep{harbor2026}, which records
for every trial the submission, the scorer output, the number of attempts and
errors, the wall-clock time, the token counts and the monetary cost. Cost and
attempts are therefore reported alongside accuracy throughout.
All backbones were called through OpenRouter, under the identifiers
\texttt{google/gemma-4-31b-it}, \texttt{google/gemma-4-26b-a4b-it} and
\texttt{z-ai/glm-4.7}, at the providers' default sampling settings; the paper
sets no temperature or context limit of its own. Costs are the per-request
figures OpenRouter reports, checked against the key balance to within $1.5\%$
on one cell; for \texttt{gemma-4-31b-it} the list price was \$0.09 per
million input and \$0.34 per million output tokens. The hub jobs ran in July and August 2026 and
the repeats on our hardware in August and September 2026, under Harbor~0.22.

Because a tool prescription presupposes the tool (Section~\ref{sec:axis}), we
state the container contents per arm. In the leaderboard arm Terminus-2 runs
on the benchmark's stock image, which ships the standard numerical and
data-frame libraries on Python~3.13. AutoDS-Tools requires the libraries its
prescription names and runs on an enriched image with them on Python~3.12.
FEDOT.LLM runs on an image built for it, with FEDOT and its dependencies on
Python~3.10. Data, splits, metrics and backbone are identical across arms; the
installed library set is not, and Section~\ref{sec:matched} measures what it is
worth.

The same applies to the instruction. Each system enters the leaderboard as it
is shipped: Terminus-2 with the task statement alone, FEDOT.LLM through a
pipeline whose tool choice is welded in, AutoDS-Tools with its
prescribed-knowledge layer applied. A row of Table~\ref{tab:leaderboard}
therefore compares deployed systems, and two further arms isolate the
architecture: AutoDS-Tools with the layer switched off, run as a full
$2\times2$ over \Ktool{} and \Kdisc{}, and Terminus-2 on the enriched image.

\subsection{Noise floor}

We measure a significance threshold and use it below. All three systems were
run three times on every TabReD task, and they do not share a floor.
Terminus-2's relative standard deviation across attempts ranges from
$\noiseTerminusMin\%$ to $\noiseTerminusMax\%$; AutoDS-Tools reproduces
itself an order of magnitude more tightly, from $\noiseAutoDSMin\%$ to
$\noiseAutoDSMax\%$, and on \allIdentAutoDS{} of its eight tasks every attempt
agrees to every digit; FEDOT.LLM agrees to every digit on \allIdentFedot{}
tasks and varies by up to $\noiseFedotMax\%$ on the rest. Between systems we
therefore treat a relative difference above $1\%$ as real, a threshold set by
the noisiest system, Terminus-2. Between configurations of AutoDS-Tools we
treat a difference above $0.1\%$ as real, a threshold set by its noisiest task.

All floors are optimistic, and for the same reason. On \bitIdentTerminus{} of
the eight tasks for Terminus-2, \bitIdentAutoDS{} for AutoDS-Tools and
\bitIdentFedot{} for FEDOT.LLM two attempts are bit-identical: the agent reproduced the same script, and on these
tasks the submission is a deterministic library call, so any two attempts that
reach the same call return the same number. The floors measure library
determinism as much as agent stability, and repeats carry less independent
information than their count suggests.

MLAgentBench has a noise floor nothing like the tabular one. Every cell was
run three times: sixty trials for AutoDS-Tools with the layer on and off and
twenty-four for Terminus-2. Three attempts of one cell
can span $\mlabOgbnOffLo$ to $\mlabOgbnOffHi$ in accuracy (AutoDS-Tools without
the layer on \texttt{ogbn-arxiv}) or $0.0844$ to $0.3846$ in micro-F1
(Terminus-2 on \texttt{fathomnet}). A difference on this benchmark is read as
real only when it exceeds a spread of that size, so Section~\ref{sec:mlab}
rests on cells in which a configuration produces no submission at all.
On TabReD a trial that exhausts its budget has written nothing; on
MLAgentBench the verifier scores whatever exists when the agent is stopped,
and such trials often submitted usable work (\texttt{clrs} scored $0.4118$
and $\mlabClrsOffHi$ on two of them), so the paper counts exhausting the
budget and returning nothing as separate events.

\subsection{Coverage}
\label{sec:coverage}

Coverage across the three systems is uneven, and every absence follows from
the architecture of FEDOT.LLM. All three systems attempt all eight TabReD
tasks; AutoDS-Tools and Terminus-2 attempt all ten MLAgentBench tasks, and
FEDOT.LLM two of the ten, because its architecture accepts tabular inputs
only: the segmentation, multi-label vision, graph and sequence tasks are
outside what it can express and would remain so after substantial
engineering. Its domain is the tabular core, the eight TabReD tasks plus
\texttt{house-price} and \texttt{spaceship-titanic}, and head-to-head
comparisons of all three systems are made there. That FEDOT.LLM can attempt
the fewest tasks follows from one decision: the pipeline is built around
FEDOT, a tabular AutoML library.

Two MLAgentBench tasks, \texttt{cifar10} and \texttt{imdb}, could not be
repeated, and the reason was measured. Sixteen attempts of sixteen exhausted
the agent budget without a submission at two concurrency settings, so
contention is not the cause. After the planning reports appear in the first
minute, a single Python process holds about eight cores continuously: a
training run. These two tasks are the heaviest to train and receive half the
budget of most others, $3600$\,s against $7200$, and their image installs the
CUDA build of PyTorch for a task that declares no accelerator, so training
falls back to the CPU. The five-epoch starter script finishes there; nothing
that improves on it does. Both tasks are omitted from every table that reports
three attempts per cell. The limit is our hardware.

\subsection{Deviations from the intended setup, and their measured cost}
\label{sec:deviations}

Two deviations from the intended design are known, and both are measured. The
first concerns the data. Earlier AutoDS-Tools runs on TabReD used a different
preparation: a $40{,}000$-row training slice spanning the official train and
validation splits, a test slice covering $16.7\%$ to $61.4\%$ of the official
test window, and an extra row-order column the official protocol withholds.
Those runs are excluded from every accuracy table in this paper; the
repair-loop counts of Section~\ref{sec:friction} come from them and are marked
as such. Table~\ref{tab:arms} uses a re-run of the same system through the
adapter that serves the full official splits, and both sets are released.
FEDOT.LLM and Terminus-2 ran through that adapter from the start. The earlier
comparison also differed in what was installed, the AutoDS-Tools arm carrying
the libraries its prescription names while the others used the stock image;
Section~\ref{sec:matched} closes that axis for Terminus-2; FEDOT.LLM runs on
its own image by construction. Its row is three attempts per task in the
system's full configuration: the best-quality preset, tuning on and a
$1200$\,s backend budget under the shared 8 cores, 16\,GB and $3600$\,s
ceiling. It reads the column types from the task schema, as the two agents
read them from the task text; with FEDOT's own type detection three of the
eight tasks returned no result. Four attempts that reached the ceiling while
a third job shared the machine were repeated at the concurrency of the rest,
and the repeats are the ones counted; every attempt is on the hub. The case
studies of Section~\ref{sec:cases} used a $2400$\,s budget under the same
ceiling.

The second concerns the image. The enriched image installs LightGBM as a dependency, but the base image and
the LightGBM wheel both lack an OpenMP runtime. The import therefore succeeds,
and the failure appears inside the fit as a broken worker pool that looks like
a modelling error. Forty-five trials met it: twenty-two of twenty-four AutoDS-Tools trials and
twenty-three of twenty-four Terminus-2 trials on the enriched image. All
forty-five diagnosed and repaired it in a single step, so it went unnoticed
until we read the traces.
Repeating the whole arm on a corrected image agrees to within $0.13\%$ on seven
tasks and $0.34\%$ on the eighth, so the defect cost no measurable accuracy. It
could still have distorted a crossed design, since the fault sits inside the
prescribed library and can strike only the cells where that prescription is
on. Every cell reported in the paper ran on the corrected image.
